\response We agree, and we have now put the two forcings on one scale. We converted the heat and freshwater flux composites of Figs. A3 and A4 into the surface density flux that each one produces, using the thermal expansion and haline contraction coefficients at the winter surface temperature and salinity of each climate state. This is the same conversion as in manuscript Figs. 3 and 4. Figure~\ref{fR:samrho} shows the result on a single colour scale, where positive values mean that the surface water becomes denser.

\rfig{0.9}{../letter_figs/figR22_sam_density_tendency_1p2sigma.png}{New analysis. Winter (JJA) high-SAM minus low-SAM composites of the surface density flux produced by (a-e) the surface heat flux and (f-j) the surface freshwater flux, sea ice included, for the five climate states, on one colour scale in kg m\textsuperscript{-2} month\textsuperscript{-1}. Positive values mean that the surface water becomes denser. The black line is the 15\% winter sea ice edge of each state. The composites use the winters in which the SAM index lies more than 1.2 standard deviations above or below its mean.\label{fR:samrho}}

On a common scale the two terms separate by ice cover. Inside the winter ice edge the freshwater term is the larger of the two over 90–93\% of the area in every state, and on average its magnitude is 6 to 9 times that of the heat term. North of the ice edge the heat term is the larger one over 72–79\% of the area, and the freshwater term is on average only 0.3 to 0.5 of the heat term. This is the same split between ice-covered and open water that governs the mean state. These are local magnitudes. The freshwater term makes the water denser at the coast and lighter further offshore, so the two terms cancel in different ways when they are summed over an area, and the comparison tells where each term acts rather than which one contributes more in total.

The corresponding statement in the Results reads:

\begin{change}
``The heat and freshwater composites are shown in their native units in Fig. A3 and Fig. A4, which makes the two forcings difficult to compare directly. Expressed as equivalent surface density tendencies, the freshwater contribution is the larger of the two wherever sea ice is present, and the heat contribution is larger only in the ice-free sector, which is the same partition that governs the mean state.''
\end{change}

\claudenote{Figure~\ref{fR:samrho} is new and could go into the Supplement next to Figs. A3 and A4. Three related points for the manuscript. Figs. A3 and A4 were made from composites at 1.0 standard deviations although their captions say 1.2; this figure uses 1.2. In the Results sentence above, ``wherever sea ice is present'' is slightly too strong (90-93\% of the ice area). Manuscript Fig. 4 and its text (line 149) give the freshwater density term about 1000 times too small, because the script drops the density of water; Fig. 3 is not affected.}
